\section{Restriction increments and coordinate selection}\label{sec:restrictions}

For a coordinate $i$, let $f_i^+$ and $f_i^-$ denote the two restrictions
of $f$, and put
\[
 \Delta_i\HH(f)=\HH(f)-\frac{\HH(f_i^+)+\HH(f_i^-)}2.
\]
The restriction framework of Gonz\'alez, MacManus and
Pereyra~\cite{GMP26} asks for a coordinate at which this entropy
increment can be charged to the influence lost in the split.
We first give an exact scalar description of the increment, then
identify several limitations of possible coordinate rules.

\subsection{The exact increment and its scalar profile}

Write $h_2(u)=-u\log u-(1-u)\log(1-u)$, with its continuous endpoint
values, and define
\begin{equation}\label{eq:kappadef}
 \kappa(u)=1+h_2(u)
 -h_2\!\left(\frac12+\sqrt{u(1-u)}\right),
 \qquad 0\le u\le1.
\end{equation}
For $S\subseteq[n]\setminus\{i\}$ set
\[
 c_S=\wh f(S),\qquad d_S=\wh f(S\cup\{i\}),\qquad
 m_S=c_S^2+d_S^2,\qquad u_S=\frac{d_S^2}{m_S},
\]
omitting terms with $m_S=0$.

\begin{proposition}[Exact restriction increment]\label{prop:splitexact}
For every Boolean $f$ and every coordinate $i$,
\begin{equation}\label{eq:splitexact}
 \Delta_i\HH(f)=\sum_{S\not\ni i}m_S\kappa(u_S).
\end{equation}
In particular, $0\le\Delta_i\HH(f)\le2$.

If $\mathcal S$ has law $P_f$, define
\[
 B_i(f)=H\bigl(\mathbf1_{\{i\in\mathcal S\}}\mid
                         \mathcal S\setminus\{i\}\bigr).
\]
Then $\Delta_i\HH(f)=B_i(f)+M_i(f)$, where
\begin{equation}\label{eq:mixingbound}
 0\le M_i(f)\le
 \frac2{\ln2}\sum_{S\not\ni i}\min(c_S^2,d_S^2)
 \le\frac2{\ln2}\I_i(f).
\end{equation}
\end{proposition}

\begin{proof}
The two restricted coefficients indexed by $S$ are $c_S+d_S$ and
$c_S-d_S$. Suppressing subscripts and putting
$v=(c+d)^2/(2m)=1/2+cd/m$, the parent pair contributes
$-m\log m+mh_2(u)$ to entropy, whereas the average contribution
of the two branches is $-m\log(2m)+mh_2(v)$.
Their difference is $m[1+h_2(u)-h_2(v)]$.
Changing the sign of $cd$ does not affect $h_2(v)$, proving
\eqref{eq:splitexact}. Both $h_2(u)$ and $1-h_2(v)$ lie in $[0,1]$,
and $\sum_Sm_S=1$.

The term $\sum_Sm_Sh_2(u_S)$ is exactly $B_i(f)$, leaving
\[
 M_i(f)=\sum_Sm_S
 \left[1-h_2\!\left(\frac12+\frac{c_Sd_S}{m_S}\right)\right].
\]
To bound it, let
\[
 g(t)=\frac{(1+t)\ln(1+t)+(1-t)\ln(1-t)}2,
 \qquad 0\le t\le1.
\]
For $0\le t<1$, the inequality
$\atanh t\le t/\sqrt{1-t^2}$ follows by comparing derivatives
from zero. Since $g'(t)=\atanh t$, integration gives
$g(t)\le1-\sqrt{1-t^2}$. Apply this with
$t=2\sqrt{u(1-u)}$, for which
$1-\sqrt{1-t^2}=2\min(u,1-u)$, and divide by $\ln2$.
Continuity handles the endpoints. Multiplication by $m_S$ and
summation yield \eqref{eq:mixingbound}.
\end{proof}

\begin{lemma}[Concavity of the scalar profile]\label{lem:kappaconcavity}
The function $\kappa$ is continuous on $[0,1]$, symmetric about
$1/2$, and strictly concave on $(0,1)$. It satisfies
$\kappa(0)=\kappa(1)=0$ and $\kappa(1/2)=2$.
Consequently,
\begin{equation}\label{eq:kappabound}
 0\le\Delta_i\HH(f)\le\kappa(\I_i(f)).
\end{equation}
As $u\downarrow0$,
\begin{equation}\label{eq:kappasmall}
 \kappa(u)=u\log(1/u)+\frac{3u}{\ln2}+O(u^2).
\end{equation}
\end{lemma}

\begin{proof}
Put $t=2\sqrt{u(1-u)}$ and
$G(t)=1-h_2((1+t)/2)$. Away from the midpoint,
\[
 h_2''(u)=-\frac4{(\ln2)t^2},\qquad
 G'(t)=\frac{\atanh t}{\ln2},\qquad
 G''(t)=\frac1{(\ln2)(1-t^2)},
\]
\[
 t'=\frac{2(1-2u)}t,\qquad t''=-\frac4{t^3}.
\]
The contributions $h_2''(u)$ and $G''(t)(t')^2$ cancel, leaving
\begin{equation}\label{eq:kappasecond}
 \kappa''(u)=-
 \frac{4\atanh(2\sqrt{u(1-u)})}
 {(\ln2)(2\sqrt{u(1-u)})^3}<0.
\end{equation}
At the midpoint, with $u=1/2+x$, the expansion is
\[
 \kappa(1/2+x)
 =2-x^2\log(1/x^2)-\frac{3x^2}{\ln2}
       +O\bigl(x^4\log(1/|x|)\bigr).
\]
Thus the first derivative tends to zero on both sides; the
concavity on the two half-intervals joins across $1/2$.
Continuity and the endpoint values follow from
\eqref{eq:kappadef}. Jensen's inequality, together with
$\sum_Sm_S=1$ and $\sum_Sm_Su_S=\I_i(f)$, proves
\eqref{eq:kappabound}. Finally,
$h_2(u)=u\log(1/u)+u/\ln2+O(u^2)$ and
$G(2\sqrt{u(1-u)})=2u/\ln2+O(u^2)$ give
\eqref{eq:kappasmall}.
\end{proof}

Since a concave function vanishing at zero has a nonincreasing
ratio to its argument, any coordinate with $\I_i(f)\ge\tau>0$
satisfies
\begin{equation}\label{eq:kappacutoff}
 \Delta_i\HH(f)\le\frac{\kappa(\tau)}{\tau}\I_i(f).
\end{equation}
For example, $\kappa(1/16)/(1/16)=8.21872678758653\ldots$.
The logarithm in \eqref{eq:kappasmall} cannot be replaced by a
uniform constant for every coordinate, as
Proposition~\ref{prop:bentsplice} below shows.

\subsection{Valid induction and block exposure}

The exact influence identity is
\begin{equation}\label{eq:restrictionI}
 \I(f)=\I_i(f)+\frac{\I(f_i^+)+\I(f_i^-)}2.
\end{equation}
Indeed, restriction averages preserve the influence of every
other coordinate, while coordinate $i$ disappears.

\begin{proposition}[A suitable-coordinate induction]\label{prop:restrictioninduction}
Suppose a class of Boolean functions is closed under restrictions
and each nonconstant member has a relevant coordinate $i$ with
$\Delta_i\HH(f)\le C\I_i(f)$. Then every member satisfies
$\HH(f)\le C\I(f)$.

Existence of a finite global FEI constant is also equivalent to
existence of a finite $K$ such that every nonconstant Boolean
function has a relevant coordinate satisfying $B_i(f)\le K\I_i(f)$.
A global FEI constant $C$ permits $K=C$; the latter erasure
inequality implies FEI with constant $K+2/\ln2$.
\end{proposition}

\begin{proof}
The first assertion follows by induction on the number of
coordinates, using \eqref{eq:restrictionI}.
For the second, the entropy chain rule gives
\begin{equation}\label{eq:hanerasure}
 \sum_i B_i(f)\le\HH(f),
\end{equation}
because conditioning on all other coordinates can only decrease
each summand of a fixed chain rule. Under $\HH(f)\le C\I(f)$,
some positive-influence coordinate must have $B_i(f)\le C\I_i(f)$.
Conversely, \eqref{eq:mixingbound} turns the proposed erasure
charge into a suitable-coordinate entropy charge, to which the
first assertion applies.
\end{proof}

The quantifier is existential at every restricted function.
Neither an arbitrary coordinate nor a preferred coordinate rule
is justified by Proposition~\ref{prop:restrictioninduction}.
The implication is hereditary: it is not the pointwise
inequality
\[
 \HH(f)\le\sum_iB_i(f)+\frac2{\ln2}\I(f),
\]
which is false for the coded examples in~\cite{BrainErasure}.
Theorem~\ref{thm:simultaneous} and Lemma~\ref{lem:plateau}
also show how finite depths of restrictions can preserve entropy
while revealing no local entropy cost.

More generally, consider a finite adaptive restriction tree, and
choose its branches by the uniform input. If $f_v$ is the function
at a visited internal node and $i(v)$ its split coordinate, then
\begin{align}
 \HH(f)&=\E\sum_{v\ {\rm internal\ on\ the\ path}}
               \Delta_{i(v)}\HH(f_v)
              +\E\HH(f_{\rm leaf}),\label{eq:treetelescopeH}\\
 \I(f)&=\E\sum_{v\ {\rm internal\ on\ the\ path}}
               \I_{i(v)}(f_v)
              +\E\I(f_{\rm leaf}).\label{eq:treetelescopeI}
\end{align}
These identities permit amortized charges and stopping at classes
with established FEI bounds, as in~\cite{GMP26}. They require the
cost along the whole path to be controlled.

\begin{proposition}[Exact block increment]\label{prop:blockincrement}
Fix a block $W\subseteq[n]$. For $T\subseteq W^c$ put
\[
 q_T=\sum_{U\subseteq W}\wh f(T\cup U)^2,\qquad
 \phi_T(y)=q_T^{-1/2}
       \sum_{U\subseteq W}\wh f(T\cup U)\chi_U(y),
\]
omitting $q_T=0$. Then
\begin{equation}\label{eq:blockincrement}
 \HH(f)-\E_y\HH(f|_{W=y})
 =H(\mathcal S\cap W\mid\mathcal S\cap W^c)
   +\sum_Tq_T\operatorname{Ent}_2(\phi_T^2),
\end{equation}
where
$\operatorname{Ent}_2 Z=\E[Z\log(Z/\E Z)]$. Moreover,
\begin{equation}\label{eq:blockmixing}
 0\le\sum_Tq_T\operatorname{Ent}_2(\phi_T^2)
 \le\frac2{\ln2}\sum_{i\in W}\I_i(f).
\end{equation}
\end{proposition}

\begin{proof}
Let $p_T(U)=\wh f(T\cup U)^2/q_T$. Then
$\E\phi_T^2=1$, and the original entropy is
$H(q)+\sum_Tq_TH(p_T)$. The restricted coefficient indexed by
$T$ is $\sqrt{q_T}\phi_T(y)$, so its entropy averaged over $y$
is
\[
 H(q)-\sum_Tq_T\operatorname{Ent}_2(\phi_T^2).
\]
Subtracting proves \eqref{eq:blockincrement}. The uniform-cube
logarithmic Sobolev inequality~\cite{OD14}, in the normalization
$D_jg=(g|_{x_j=1}-g|_{x_j=-1})/2$, states
\[
 \operatorname{Ent}_2(g^2)
 \le\frac2{\ln2}\sum_j\E(D_jg)^2.
\]
Applying it to each $\phi_T$ and using Parseval gives
\[
 \sum_Tq_T\sum_{i\in W}\E(D_i\phi_T)^2
 =\sum_{T,U}|U|\wh f(T\cup U)^2
 =\sum_{i\in W}\I_i(f),
\]
which proves \eqref{eq:blockmixing}.
\end{proof}

Block exposure therefore retains a conditional erasure charge.
Taking $W=[n]$ makes that charge $\HH(f)$ itself. Small erasure
for bounded blocks at the initial function does not replace the
hereditary requirement, because the separated supports in
Lemma~\ref{lem:plateau} have zero entropy increment under all
sufficiently shallow restrictions.

\subsection{Character-centered budgets}

Define
\[
 \D(f)=\sum_j t(\I_j(f)),\qquad t(u)=\min(u,1-u).
\]
Character multiplication preserves entropy and complements the
influences on its selected coordinates. Hence
$\D(f)=\min_U\I(\chi_Uf)$ and $\D$ is character invariant.
For a split define
\[
 d_i^{\D}(f)=\D(f)
            -\frac{\D(f_i^+)+\D(f_i^-)}2.
\]

\begin{proposition}[Centered budget and its collapse]\label{prop:centerbudget}
For every $f,i$,
\begin{align}
 d_i^{\D}(f)
 ={}&t(\I_i(f))\notag\\
 &+\sum_{j\ne i}
 \left[t(\I_j(f))
       -\frac{t(\I_j(f_i^+))+t(\I_j(f_i^-))}2\right]
 \ge t(\I_i(f)).\label{eq:centerbudget}
\end{align}
Both $d_i^{\D}$ and $\Delta_i\HH$ are character invariant.
If $\max_j\I_j(f)\le1/4$, then
\begin{equation}\label{eq:centerbudgetcollapse}
 d_i^{\D}(f)=\I_i(f)\qquad\text{for every }i.
\end{equation}
\end{proposition}

\begin{proof}
For $j\ne i$, the two child influences average to $\I_j(f)$.
Concavity of $t$ proves \eqref{eq:centerbudget}.
Restriction commutes with character multiplication up to a
character and an output sign, proving the invariances.
Each child influence is at most twice its parent's influence.
Under the stated upper bound, all parent and child influences
are at most $1/2$, so $\D=\I$ at these three nodes.
Equation~\eqref{eq:restrictionI} now gives
\eqref{eq:centerbudgetcollapse}.
\end{proof}

A suitable-coordinate inequality
$\Delta_i\HH\le C d_i^{\D}$ would telescope to
$\HH\le C\D$, stopping at characters, whose entropy is zero.
The centered budget is sometimes larger than the original
budget: for a two-input conjunction it reduces the local ratio
from four to two. It supplies no improvement in
\eqref{eq:centerbudgetcollapse}. More generally, if
$\max_j\I_j(f)\le2^{-(r+1)}$, every restriction fixing at most
$r$ coordinates has influences at most $1/2$, so the first
$r$ levels have no centering gain. This follows by iterating
the factor-two bound. The low-influence reductions
in~\cite{BrainCritical} and Theorem~\ref{thm:simultaneous}
therefore remain applicable.

\subsection{Bad coordinates with arbitrarily small influences}

The following bent-splice mechanism also underlies the local
allocation obstruction in~\cite{BrainAllocation}.

\begin{proposition}[Bent splice]\label{prop:bentsplice}
For each even $m\ge4$, there is a Boolean function $F$ on $m+1$
coordinates, with distinguished coordinate $z$, such that
\[
 \I_z(F)=2^{-m},\qquad
 \I_j(F)=\frac12+2^{-m}\quad(1\le j\le m),
\]
and
\begin{equation}\label{eq:bentspliceratio}
 \frac{\Delta_z\HH(F)}{\I_z(F)}
 \ge\log\!\left(2^m-2\cdot2^{m/2}+2\right)\ge m-1.
\end{equation}
\end{proposition}

\begin{proof}
Use binary coordinates for the $m$ base variables, with characters
$(-1)^{S\cdot x}$, and put
\[
 g(x)=(-1)^{x_1x_2+x_3x_4+\cdots+x_{m-1}x_m},\qquad
 h(x)=g(x)-2\mathbf1_{\{x=0\}}.
\]
Both are Boolean. Define $F(x,1)=g(x)$ and $F(x,-1)=h(x)$.
Set $N=2^m$, $\alpha=N^{-1}$ and $R=\sqrt N$.
The two-bit quadratic factor has four Fourier coefficients of
magnitude $1/2$; tensorization gives
$\wh g(S)\in\{-1/R,1/R\}$ for every $S$.
Each base influence of $g$ is $1/2$.
Since $g(0)=g(e_j)=1$, flipping the value at zero changes the
edge $\{0,e_j\}$ from agreeing to disagreeing, increasing
$\I_j(h)$ by $2/N$. Averaging the branches proves the displayed
influence identities.

For the $z$-split, the coefficient pairs are
$c_S=\wh g(S)-\alpha$ and $d_S=\alpha$. Thus
\[
 \frac{m_S}{d_S^2}=(\pm R-1)^2+1
                   \ge N-2\sqrt N+2.
\]
The $d_S^2\log(m_S/d_S^2)$ term alone gives
\[
 \Delta_z\HH(F)\ge B_z(F)
 \ge N\alpha^2\log(N-2\sqrt N+2).
\]
Since $\I_z(F)=\alpha$, the first inequality in
\eqref{eq:bentspliceratio} follows. The second follows from
$N-2\sqrt N+2-N/2=(\sqrt N-2)^2/2\ge0$.
\end{proof}

\begin{corollary}[Failure of two maximum-budget rules]\label{cor:badmaximum}
There are Boolean functions $\widetilde F_m$ with
$\max_i\I_i(\widetilde F_m)\le2^{-m}$ for which the unique
maximum-influence coordinate has local ratio at least $m-1$.
It is also the unique coordinate maximizing $d_i^{\D}$, with
the same ratio $\Delta_i\HH/d_i^{\D}$.
\end{corollary}

\begin{proof}
In Proposition~\ref{prop:bentsplice}, replace each base coordinate
by a majority on a disjoint block of
$K=2N^2+1$ bits, keeping $z$ unchanged. Its leaf influence is
\[
 \beta_K=\binom{K-1}{(K-1)/2}2^{-(K-1)}
 \le\frac1{\sqrt{(K+1)/2}}<\alpha.
\]
For completeness, the bound
$\binom{2t}{t}/4^t\le1/\sqrt{t+1}$ follows by induction:
the consecutive ratio is $(2t+1)/(2t+2)$, and the squared
comparison reduces to $3t+2\ge0$.

Lemma~\ref{lem:balanced} gives each new leaf influence
$(1/2+\alpha)\beta_K<\alpha$ and preserves $\I_z=\alpha$.
It also preserves $\Delta_z\HH$: the additional entropy
in the parent is $\sum_j\I_j(F)\HH(\Maj_K)$, and the
additional branch entropies average to the same quantity
by \eqref{eq:restrictionI} applied to the other coordinates.
Thus $z$ is the unique maximum-influence coordinate, and
\eqref{eq:bentspliceratio} persists.
All influences are at most $\alpha\le1/16$, so
Proposition~\ref{prop:centerbudget} gives
$d_i^{\D}=\I_i$ at every coordinate, proving the second assertion.
\end{proof}

These examples reject the two specified rules, not the
existential hypothesis of Proposition~\ref{prop:restrictioninduction}.
The cost may still be paid by choosing a different coordinate
or by an amortized argument.

\subsection{Separable influence potentials}

\begin{proposition}[An obstruction to charging every split]\label{prop:separablepotential}
Let $\phi:[0,1]\to\mathbb R$ satisfy $\phi(0)=0$,
$\phi(t)=O(t)$ as $t\downarrow0$, and have a finite left derivative
at $1/2$. Put $P_\phi(f)=\sum_j\phi(\I_j(f))$.
There is no finite $C>0$ such that
\begin{equation}\label{eq:falsepotentialcharge}
 \Delta_i\HH(f)\le C\left[
 P_\phi(f)-\frac{P_\phi(f_i^+)+P_\phi(f_i^-)}2\right]
\end{equation}
for every Boolean $f$ and every coordinate $i$.
The assertion remains false after adding to $P_\phi(f)$ any
real-valued function of $\E f$.
\end{proposition}

\begin{proof}
Before flipping zero in Proposition~\ref{prop:bentsplice},
multiply $g$ by the character on all its base coordinates.
Its spectrum remains flat, and its base influences remain
$1/2$. Now the edges from zero to $e_j$ disagree, so flipping
zero makes the other branch's influences $1/2-2\alpha$ and
the parent influences $1/2-\alpha$.
The selector still has influence $\alpha$, and the lower bound
$\Delta_z\HH\ge(m-1)\alpha$ is unchanged.
The potential drop is exactly
\[
 \phi(\alpha)+m\left[
 \phi(1/2-\alpha)
 -\frac{\phi(1/2)+\phi(1/2-2\alpha)}2\right].
\]
The bracket is $o(\alpha)$ by the one-sided derivative
assumption. Dividing the displayed expression by $m\alpha$
therefore gives a quantity tending to zero, whereas
$\Delta_z\HH/(m\alpha)\ge1-1/m$.
This contradicts \eqref{eq:falsepotentialcharge}.

Finally multiply the entire splice by a fresh independent sign.
Parent and both selector branches then have mean zero, so any
mean-only correction cancels. The new coordinate has influence
one at all three nodes, and its contribution to the potential
drop also cancels. Entropy and the selector increment are
unchanged.
\end{proof}

The assumptions include $\phi(t)=t$ and
$\phi(t)=\min(t,1-t)$. The proposition concerns a charge at
every coordinate; it does not rule out a suitable-coordinate
inequality, a nonseparable potential, or a global tree budget.

\subsection{Finite checks and exact witnesses}\label{subsec:restrictionchecks}

Enumeration of all $2^{16}=65,536$ Boolean truth tables on four
labelled coordinates, including functions with unused inputs,
gives the values in Table~\ref{tab:restrictionchecks}.
Walsh numerators, influences and centered budgets are calculated
with exact integer or dyadic arithmetic. The ordering of
entropy ratios in the enumeration uses floating-point
logarithms; the exact witness formulas below were also evaluated
with 80-digit arithmetic. Thus the table reports numerical
finite maxima, not interval-certified comparisons or
dimension-independent bounds.

\begin{table}[ht]
\centering
\caption{Observed maxima among four-variable Boolean functions.
Ratios with zero denominator are omitted.}\label{tab:restrictionchecks}
\begin{tabular}{lc}
\toprule
Quantity & Observed maximum\\
\midrule
$\HH(f)/\I(f)$ & $3.402475551198587$\\
$\min_{i:\I_i>0}\Delta_i\HH(f)/\I_i(f)$ & $4.867668746754797$\\
$\min_{i:d_i^{\D}>0}\Delta_i\HH(f)/d_i^{\D}(f)$
 & $4.867668746754797$\\
$\max_{i:\I_i>0}\Delta_i\HH(f)/\I_i(f)$
 & $5.921783780331384$\\
$\max_{i:d_i^{\D}>0}\Delta_i\HH(f)/d_i^{\D}(f)$
 & $5.921783780331384$\\
\bottomrule
\end{tabular}
\end{table}

The entropy-to-influence witness is a four-input conjunction,
up to output sign, with exact ratio
\[
 12-\frac{49}{16}\log 7.
\]
The suitable-coordinate witness is a three-input conjunction
and one unused variable. Its squared coefficients are $9/16$
at the empty set and $1/16$ at each of seven nonempty sets;
each relevant influence and centered budget is $1/4$.
Its local ratio is exactly
\[
 12-\frac92\log 3
 =4.86766874675479718345817475223482571058\ldots .
\]
For an arbitrary-coordinate witness, take the function on
$\{0,1\}^4$ equal to $-1$ precisely at
$0000,0001,0010$. Its squared coefficients are $25/64$ once,
$9/64$ three times, and $1/64$ twelve times. At either of its
coordinates of influence $1/8$, the local ratio is
\[
 24-\frac{25}{4}\log 5-\frac94\log 3
 =5.92178378033138391753959094180872425614\ldots .
\]
The exact formulas establish these witness values independently
of the numerical ordering. Larger self-composition examples
in~\cite{BrainCritical} require larger suitable-coordinate
constants, so the finite table cannot be used as a uniform
restriction bound.
